% Transcription — fonds Alexandre Grothendieck, shelfmark 1, batch 4 (pages 61-80).
% Established from the facsimile archives/batches/1/batch-04.pdf (a mirror of
% https://grothendieck.umontpellier.fr/1.pdf), per the skill
% transcribe-grothendieck.
%
% Pass: Fable 5.1 (claude-fable-5-1), 2026-09-10 — first pass, unchecked against
% the pages by a human.
%
% All twenty pages are transcribed. \page{} numbers the position in the
% folder, as batches 1 to 3 did: the archivists' pencil numbers run two
% ahead, so page 61 here is pencilled 63 and page 80 is pencilled 82.
%
% Pages 61 and 62 are one sheet photographed twice. A slip is pasted to its
% upper left; page 61 shows it folded back (the sheet's own text visible
% underneath), page 62 shows it unfolded over the body. The slip carries the
% archivists' number 64 and is transcribed under \page{62}; the text it
% covers is transcribed under \page{61}. The slip's text — the orthonormal
% system e_J — belongs, by its formula number (5), between the Proposition
% of page 61 and its « En particulier ».
%
% The batch falls into four runs.
%
%   Pages 61-69 — the continuation of « 3. Produits extérieurs hilbertiens »
%     begun at page 58 of batch 3, in the same fast blue-black hand on
%     yellowed sheets: (4) (a₁∧…∧a_n, b₁∧…∧b_n) = det((a_i,b_j)), the
%     Proposition on the Hilbert structure of ΛE, (5) ‖a₁∧…∧a_n‖ =
%     det((a_i,a_j))^{1/2} ≤ ∏‖a_i‖ and (6) Hadamard's inequality
%     |det A| ≤ ∏‖a_i‖, Corollaire 1 |det((a_i,b_j))| ≤ ∏‖a_i‖‖b_i‖; then
%     u₁∧…∧u_n for operators, (8) its factorisation through the
%     antisymmetriser, (9) ‖u₁∧…∧u_n‖ ≤ ∏‖u_i‖, (10) (Λv)(Λu) = Λ(vu),
%     (11) (Λu)* = Λ(u*), (12) (Λh)^α = Λ(h^α) for h ≥ 0 with its proof
%     through rational α; Proposition 1 summing these up; Proposition 2
%     (the eigenvalues of Λ^n u are the products λ_{i₁}…λ_{i_n}) with a
%     proof through the decomposition ΛE = Σ Λ^p E₀ ⊗ Λ^q E₁; Corollaire 1
%     with (21) λ₁…λ_n = λ₁(Λ^n u), (22) ρ₁…ρ_n = ‖Λ^n u‖ rewritten in the
%     margin of page 68, (23) |λ₁(u)…λ_n(u)| ≤ ρ₁(u)…ρ_n(u), and (24)
%     ρ₁(uv)…ρ_n(uv) ≤ ∏ρ_i(u)ρ_i(v), attributed to H. Weyl.
%   Pages 70-73 — a fair copy on whiter sheets, torn along the top, headed
%     « 3. Puissances extérieures d'op. de Fredholm, inégalités diverses »,
%     with its own numbering restarting at (1): ‖u₁∧…∧u_n‖₁ ≤ (1/n!)∏‖u_i‖₁,
%     α_n(u) = Tr Λ^n u, |α_n(u)| ≤ α_n(|u|) = ‖Λ^n u‖₁ ≤ (1/n!)‖u‖₁^n,
%     |det(1+u)| ≤ det(1+|u|), Théorème 1 det(1+zu) = ∏(1+zλ_i(u)) of
%     genus 0, (11) the bound on (Λ^p u)∧(Λ^q v), (14) α_n(|u+v|) ≤
%     Σ α_p(|u|)α_{n-p}(|v|), Théorème 2 (15)-(17) on σ_n^m(u+v),
%     Prop. 3 (19) σ_n^m(u) = sup |α_m(v_n u)| over ‖v_n‖ ≤ 1 of rank ≤ n,
%     Prop. 4 (18) ∏(1+ρ_i(u)) = sup det(1+v_n u).
%   Pages 74-75 — two brownish sheets in another pen: the subadditivity of
%     Σ f_φ(ρ_i(w)) for f_φ(ρ) = ∫φ(r)/(1/ρ+r) dr, integration by parts
%     against a decreasing φ, the case φ(r) = r^{-p} giving S_p(u+v) ≤
%     S_p(u) + S_p(v) with S_p(u) = Σ ρ_i(u)^p, and a « Question » on the
%     functions f for which Σ f(ρ_i(·)) is subadditive.
%   Pages 76-79 — an earlier draft of the fair copy, on brownish sheets,
%     headed in red pencil « Mettre 1 : N° 2 et cor. N° — »: Proposition 3
%     (13), Corollaire 1 (14), Corollaires 2-5 (15)-(19), Théorème 2 (20)
%     det(1+u) = ∏(1+λ_i), (21) det(1+zu) = ∏(1+zλ_i), and its proof, the
%     lower half of page 78 and all of page 79 crossed out diagonally.
%   Page 80 — a note on a convex function f on an open set U: the epigraph
%     G_f, a supporting hyperplane at (a, f(a)), the affine minorant L_a,
%     f(x) = sup_a L_a(x).
%
% Cancellations: pages 61, 62, 64, 66, 67, 68, 69, 70, 71, 72, 73, 76, 77,
% 78, 79 and 80 carry lines or blocks struck through; each is transcribed
% as \struck{} with one \note{} for the block. Red pencil frames formulas
% (4), (5), (10), (11), (12), (23), (24) of the first run, (13), (14), (15),
% (18), (20), (23) of the draft, and underlines several words of the fair
% copy. The long sideways notes in the left margins of pages 61, 66, 68,
% 71, 78 and 79 are recorded as \marginal{} where they can be read; the
% pencil note of page 61 and the struck note of page 66 are largely lost.
%
% Notation: he writes Λ with « n » and « (2) » above, ⊗ with « (2) » above;
% both are set as superscripts, \Lambda^{n\,(2)} and \otimes^{(2)}, as in
% batch 3. Λ alone, without exponent, is kept where he writes it so: he
% shifts between Λu and Λ^n u on the same page. The Fredholm class is L¹
% or L^{(1)} by turns and is set L^{1}; ‖·‖₁ and ‖·‖₂ are the trace and
% Hilbert-Schmidt norms of batch 3; ρ_i(u) the singular values, λ_i(u)
% the eigenvalues, α_n(u) the coefficients of det(1+zu) of « Chap. 1 »,
% |u| = (u*u)^{1/2}. S_n is his gothic S for the symmetric group.
%
% The dating is the inventory's own for the shelfmark. Its group,
% « MATHÉMATIQUES / RECHERCHE / Travaux », carries no date.
\documentclass[11pt,a4paper]{article}
% Le préambule commun à toutes les transcriptions du fonds Grothendieck.
%
% Il tient en une idée : **l'incertitude se compose, elle ne se résout pas.**
% Une transcription qui lisse un mot douteux a détruit la seule chose pour
% laquelle on la faisait — un lecteur doit pouvoir savoir, à chaque endroit,
% ce qui était sur le feuillet et ce qui a été ajouté. Les six macros qui
% suivent sont donc l'appareil critique, et non de la décoration.
%
% Les mêmes commandes sont reconnues par `scripts/render.mjs`, qui produit la
% vue de lecture du site. Une macro ajoutée ici doit l'être aussi là-bas,
% faute de quoi le rendu échouera bruyamment — ce qui est voulu : un rendu
% silencieusement faux est pire qu'un rendu absent.

\ProvidesPackage{grothendieck}[2026/08/08 Transcriptions du fonds Grothendieck]

\usepackage{amsmath,amssymb,amsthm}
\usepackage{mathrsfs}
\usepackage{stmaryrd}
% Les diagrammes commutatifs sont partout dans ce fonds : c'est la langue
% même de la géométrie algébrique de Grothendieck, et une transcription qui
% les rendrait en prose ne transcrirait rien.
\usepackage{tikz-cd}
% Français par défaut — c'est la langue des feuillets — mais l'anglais est
% chargé aussi : la traduction et le résumé sont des documents anglais, et la
% typographie française (espaces avant les deux-points) y serait une faute.
% Ces documents basculent par \selectlanguage{english} en tête de corps.
\usepackage[english,french]{babel}
\usepackage{fontspec}
\usepackage{geometry}
\geometry{a4paper,margin=25mm,marginparwidth=28mm}
\usepackage{marginnote}
\usepackage{xcolor}
% ulem, not soul. `soul`'s \st and \ul are built on a character-by-character
% scan that XeLaTeX's font machinery breaks: the document fails with an
% undefined control sequence rather than degrading. ulem does the same two jobs
% and is Unicode-engine safe. `normalem` stops it hijacking \emph, which the
% transcriptions use for Grothendieck's own emphasis.
\usepackage[normalem]{ulem}
\usepackage{hyperref}
\hypersetup{colorlinks=true,linkcolor=black,urlcolor=black,pdfborder={0 0 0}}

% --- Métadonnées du lot -----------------------------------------------------
% Reprises telles quelles par le script de rendu, qui les affiche en tête de
% la vue de lecture. Elles fixent ce que le fichier prétend couvrir : sans
% elles, une transcription flotte sans point d'attache dans un fonds de seize
% mille feuillets.

\newcommand{\folder}[1]{\gdef\grothFolder{#1}}
\newcommand{\batch}[1]{\gdef\grothBatch{#1}}
\newcommand{\leaves}[2]{\gdef\grothFirst{#1}\gdef\grothLast{#2}}
\newcommand{\foldertitle}[1]{\gdef\grothFoldertitle{#1}}
\gdef\grothFolder{?}\gdef\grothBatch{?}\gdef\grothFirst{?}\gdef\grothLast{?}\gdef\grothFoldertitle{}

% --- Le résumé --------------------------------------------------------------
%
% Il ouvre la lecture modernisée, sous le titre « Résumé », et il est composé
% en petit. Ce n'est pas un ornement : c'est le seul passage du document qui
% ne soit pas des mathématiques, et un lecteur doit pouvoir le reconnaître
% avant de l'avoir lu — puis le sauter s'il connaît déjà le sujet, ou s'y
% arrêter s'il ne le connaît pas. Le filet qui le suit marque la reprise du
% corps.

\newenvironment{resume}
  {\small\setlength{\parskip}{0.45em}}
  {\par\medskip\hrule\medskip}

% --- Mots-clés ---------------------------------------------------------------
%
% En anglais, en clôture du résumé de la lecture modernisée : le vocabulaire
% moderne sous lequel chercher ce que les feuillets construisent. Le manifeste
% du site (`npm run manifest`) les extrait pour étiqueter la cote entière —
% cette ligne est la source unique des tags, il n'y a pas de fichier à côté.
\newcommand{\keywords}[1]{\par\smallskip\noindent
  {\footnotesize\textsc{Keywords} --- \textit{#1}}\par}

% --- L'appareil critique ----------------------------------------------------

% Le numéro de feuillet, en marge. C'est l'ancre qui permet de régler un
% désaccord sur une formule en regardant *un* feuillet plutôt qu'une chemise
% de six cents. Le site s'en sert aussi pour tourner les pages du fac-similé
% quand on fait défiler la transcription.
\newcommand{\leaf}[1]{%
  \marginnote{\footnotesize\color{gray!70}\textsf{#1}}%
  \label{leaf:#1}\ignorespaces}

% Illisible : on ne devine pas. Un mot inventé qui se lit comme les autres est
% le pire résultat possible.
\newcommand{\ill}{\textcolor{red!60!black}{[\,\ldots\,]}}

% Lecture douteuse : le mot est donné, et signalé comme incertain.
\newcommand{\uncertain}[1]{\uline{#1}}

% Ajout de l'éditeur — une lettre manquante, un mot sous-entendu.
\newcommand{\add}[1]{\textcolor{blue!50!black}{[#1]}}

% Biffé par Grothendieck lui-même. Ce qu'il a rayé fait partie du document :
% une rature dit souvent où la pensée a changé de direction.
\newcommand{\struck}[1]{\sout{#1}}

% Note du transcripteur, sur l'état matériel du feuillet.
\newcommand{\note}[1]{\par\smallskip{\small\color{gray!60!black}\textit{[#1]}}\par\smallskip}

% Note marginale de Grothendieck, distincte du corps du texte.
\newcommand{\marginal}[1]{\par\smallskip\noindent\hspace{1em}%
  {\small\color{gray!60!black}#1}\par\smallskip}

% --- Titre ------------------------------------------------------------------

\AtBeginDocument{%
  \begin{center}
    {\large\bfseries Fonds Alexandre Grothendieck --- cote n\textsuperscript{o}~\grothFolder}\\[2pt]
    {\itshape\grothFoldertitle}\\[4pt]
    {\small feuillets \grothFirst--\grothLast\ (lot \grothBatch)}\\[2pt]
    {\footnotesize\color{gray!60!black}Transcription établie d'après le fac-similé numérisé
     par l'Université de Montpellier.\\
     Les lectures incertaines sont soulignées, les illisibles notées [\,\ldots\,],
     les ajouts entre crochets.}
  \end{center}
  \vspace{1em}}

\endinput


\folder{1}
\batch{4}
\pages{61}{80}
\dating{1953}
\watermark{Édition de démonstration}
\foldertitle{Analyse fonctionnelle : notes manuscrites (s.d.), lettre (1953)}

\begin{document}

\section{Produits extérieurs hilbertiens (suite)}

\note{Titre de l'éditeur. Les pages 61 à 69 poursuivent le « 3. Produits extérieurs hilbertiens. Applications : des inégalités diverses » ouvert page 58 du lot 3 ; la numérotation des formules continue celle de ce paragraphe, dont la page 60 s'arrêtait sur la formule $(a_1 \wedge \cdots \wedge a_n, b_1 \wedge \cdots \wedge b_n) = \det((a_i,b_j))$ et une parenthèse ouverte.}

\page{61}
\[
  \text{(4)}\qquad (a_1 \wedge \cdots \wedge a_n,\; b_1 \wedge \cdots \wedge b_n) = \det\bigl((a_i, b_j)\bigr)_{i,j \leq n}
\]
\note{formule encadrée au crayon rouge}
\uncertain{On} \uncertain{peut} \uncertain{donc} \uncertain{énoncer} \uncertain{les} \uncertain{premières} \uncertain{propriétés} \ill{}

\emph{Proposition.} Soit \struck{un} $E$ un espace de Hilbert. Il existe sur $(\widehat{\Lambda} E) \times (\widehat{\Lambda} E)$ \note{le $\Lambda$ porte un accent circonflexe ; ainsi sur la page} une forme hermitienne définie positive (et une seule), satisfaisant à (4). \uncertain{La} \uncertain{complétée} \uncertain{est} \uncertain{le} \struck{produit tensoriel} \uncertain{hilbertien} \ill{} \uncertain{de} $\widehat{\Lambda} E$. Si $(e_i)_{i \in I}$ est une base \uncertain{hilbertienne} de $E$, alors les $e_J$ définis par (5), forment une base \struck{orthonormale} \uncertain{hilbertienne} \struck{\ill{}} \note{deux lignes lourdement biffées} dans $\Lambda^{(2)} E$. \note{les $e_J$ et la formule (5) sont définis sur le feuillet collé, page 62} En particulier on a
\[
  \text{(5)}\qquad \|a_1 \wedge \cdots \wedge a_n\| = \bigl(\det((a_i,a_j))\bigr)^{1/2} \leq \|a_1\| \cdots \|a_n\| \qquad (a_i \in E)
\]
\note{formule encadrée au crayon rouge ; le numéro « (5) » y est marqué d'un signe, peut-être « (5') », le feuillet collé portant déjà un (5) ; sous la formule, « $a_i$ » et « (5) » répétés}
En effet, l'égalité est un cas particulier de (4) \ill{}. \struck{\ill{}}, \uncertain{donc} \uncertain{l'expression} \ill{} \uncertain{les} \uncertain{bases} \uncertain{orthonormales} $(e_i)$, \struck{\ill{}} posons $a_i = \sum_j a_{ij}\, e_j$ \note{le $i$ est cerclé} --- \struck{Vérifier aussitôt que} soit $A = (a_{ij})_{i,j \leq n}$, alors $(a_i, a_j) = \sum_k a_{ik}\, \bar{a}_{jk}$ \struck{$\det((a_i,a_j)) = \det$} est l'élément d'indice $i,j$ de $A A^{*}$, donc
\[
  \det\bigl((a_i,a_j)\bigr) = \det A A^{*} = \det A \cdot \det A^{*} = \det A \cdot \overline{\det A} = |\det A|^{2},
\]
d'où \uncertain{il} \ill{}
\[
  \text{(6)}\qquad |\det A| \leq \|a_1\| \cdots \|a_n\| \qquad \text{(inégalité de Hadamard)},
\]
\uncertain{puisque} \ill{} \uncertain{pour} \uncertain{la} \uncertain{calculer}, \ill{} \note{trois lignes de prose, en grande partie illisibles : il s'agit du choix des $e_i$} \uncertain{il} \uncertain{suffit} \uncertain{de} \ill{}. \uncertain{Pour} \uncertain{ceci}, \ill{} \uncertain{vecteurs} \ill{} les $(e_i)$ \uncertain{choisis} \uncertain{de} \uncertain{façon} \uncertain{que} $A$ \uncertain{soit} \struck{\ill{}} \uncertain{triangulaire}, d'où $\det A = a_{11}\, a_{22} \cdots a_{nn}$, ce qui achève la démonstration \ill{}.

\emph{Corollaire 1.} Soient $a_i, b_i$ ($i = 1, \ldots, n$) des éléments d'un espace de Hilbert $E$. Alors on a
\[
  \bigl|\det\bigl((a_i, b_j)\bigr)\bigr| \leq \prod_i \|a_i\|\, \|b_i\| .
\]
\marginal{\ill{}} \note{longue note au crayon dans la marge gauche, avec des flèches, illisible même redressée}

\page{62}
\note{Feuillet collé sur la page 61, numéroté 64 par les archivistes ; il est ici déplié et couvre le milieu de la page. Sa moitié inférieure est biffée.}
Soit $(e_i)_{i \in I}$ une base hilbertienne de $E$ \note{en interligne : « $I$ \uncertain{tot.} \uncertain{ordonné} »}, pour toute partie \struck{finie} $J = (i_1, \ldots, i_n)$, $i_1 < i_2 < \cdots < i_n$ de $I$, posons
\[
  \text{(5)}\qquad e_J = e_{i_1} \wedge \cdots \wedge e_{i_n} .
\]
\note{formule encadrée}
On voit aussitôt sur (4) que $(e_J, e_{J'}) = \delta_{JJ'}$ (symbole de Kronecker), donc les $e_J$ forment un système orthonormal dans $\Lambda^{(2)} E$ \note{un mot en interligne, illisible}. \struck{\ill{} \uncertain{le} \uncertain{sous-espace} \ill{} \uncertain{engendré} \ill{} \uncertain{est} \uncertain{évidemment} \uncertain{dense}, \ill{} \uncertain{de} $E$, \ill{} \uncertain{dans} $\Lambda E$ \ill{} \uncertain{l'application} \uncertain{multilinéaire}} \note{deux lignes biffées} \uncertain{on} \ill{} \uncertain{les} $e_J$ \ill{} \note{en interligne : « dans $\Lambda^{(2)} E$ »} \uncertain{identifier} \ill{} \uncertain{engendré} \uncertain{par} \uncertain{l'image} \uncertain{de} $(E)^{n}$ \uncertain{par} \uncertain{l'application} \uncertain{multilinéaire} \uncertain{naturelle} \uncertain{de} $E^{n}$ \uncertain{dans} $\Lambda E$, \uncertain{donc} \uncertain{dense} \uncertain{dans} $\Lambda E$ \uncertain{puisque} \ill{} $E_0$ \ill{} \uncertain{dense} \uncertain{dans} $E$, \uncertain{donc} \uncertain{dense} \uncertain{dans} $\Lambda^{(2)} E$.

\page{63}
En effet, il suffit d'utiliser (4), de l'inégalité de \ill{} \uncertain{appliquée} \uncertain{à} \ill{} \uncertain{Cauchy-Schwarz} \uncertain{pour} \uncertain{la} \uncertain{première} \ill{}, et d'utiliser (5).

Ce corollaire \uncertain{prouve} que la \uncertain{constante} $K_n$ \uncertain{qui} \uncertain{intervenait} dans (Chap.~1, \S 2, N° 2, prop 1) est $\leq 1$ pour l'espace de Hilbert, (compte tenu de l'identification du dual $E'$ de $E$ à $\bar{E}$).

Soient maintenant $E$, $F$ deux espaces de Hilbert, et soient $u_1, \ldots, u_n$ des \uncertain{opérateurs} linéaires continus de $E$ dans $F$, considérons leur produit \uncertain{extérieur} $u_1 \wedge \cdots \wedge u_n$ (Chap.~1, \S 1, N° 1). On a :
\[
  \text{(8)}\qquad u_1 \wedge \cdots \wedge u_n = n!\; \varphi_n \cdot (u_1 \otimes \cdots \otimes u_n) \cdot a_n
\]
\note{la formule (7) n'apparaît pas sur ces pages}
où $a_n$ est \uncertain{l'antisymétriseur} \struck{\ill{}} \note{un mot en interligne} de $\Lambda E$ dans $\bigotimes E$, et $\varphi_n$ l'homomorphisme naturel de $\bigotimes F$ \uncertain{sur} $\Lambda F$. Comme $\frac{1}{\sqrt{n!}}\, a_n$ \ill{} est une \uncertain{isométrie} \ill{} \uncertain{préhilbertiens}, et $\sqrt{n!}\; \varphi_n$ une \ill{} de norme $\leq 1$ pour les structures hilbertiennes \ill{} $\bigotimes F$ et $\Lambda F$, on voit que
\[
  \|u_1 \wedge \cdots \wedge u_n\| = \Bigl\|\sqrt{n!}\; \varphi_n \cdot (u_1 \otimes \cdots \otimes u_n) \cdot \tfrac{1}{\sqrt{n!}}\, a_n\Bigr\| \leq \|u_1 \otimes \cdots \otimes u_n\|,
\]
d'où en vertu de (2) :
\[
  \text{(9)}\qquad \|u_1 \wedge \cdots \wedge u_n\| \leq \|u_1\| \cdots \|u_n\| .
\]
Par suite, $u_1 \wedge \cdots \wedge u_n$ peut être étendu en une application linéaire continue de $\Lambda^{(2)} E$ dans $\Lambda^{(2)} F$ \note{en interligne : « notée encore $u_1 \wedge \cdots \wedge u_n$ »}, dont la norme satisfait à (9). Bien entendu, les règles de calcul usuelles subsistent pour les \note{en interligne : « (quelques uns (extensions)) »} \struck{\ill{}} $u_1 \wedge \cdots \wedge u_n$. \uncertain{Alors} \ill{} \uncertain{particulières}, \uncertain{si} \uncertain{on} \uncertain{a}
\note{la phrase se poursuit page 64}

\page{64}
\ill{} $u \in L(E,F)$, $v \in L(F,G)$, alors on aura
\[
  \text{(10)}\qquad (\Lambda v)(\Lambda u) = \Lambda(vu)
\]
\note{le numéro (10) est cerclé au crayon rouge}
(Notons aussi la formule
\[
  \text{(11)}\qquad \struck{(\Lambda u)^{*} = \ldots}\qquad (u_1 \wedge \cdots \wedge u_n)^{*} = u_1^{*} \wedge \cdots \wedge u_n^{*}
\]
\note{le numéro (11) cerclé au crayon rouge ; une accolade réunit cette formule à « $(\Lambda u)^{*} = \Lambda(u^{*})$ » écrit en interligne}
Pour l'établir, on peut \ill{} \struck{\ill{}} \note{trois lignes biffées} \ill{} formule (ou $u_1 \wedge \cdots \wedge u_n = \frac{1}{n!} \struck{\ill{}} \sum_{0 \leq p \leq n,\; i_1 < \cdots < i_p \leq n} (-1)^{\ill{}}\, \Lambda(u_{i_1} \cdots u_{i_p})$ \note{lecture incertaine de la formule interne ; il s'agit d'un développement du produit extérieur} --- \uncertain{voir} \uncertain{Introduction}, 2, formule (8), ---), on a \ill{}
\begin{gather}
  \bigl((\Lambda u)\, a_1 \wedge \cdots \wedge a_n,\; b_1 \wedge \cdots \wedge b_n\bigr) = \det\bigl((u a_i, b_j)\bigr) = \det\bigl((a_i, u^{*} b_j)\bigr) \\
  = \bigl(a_1 \wedge \cdots \wedge a_n,\; (\Lambda u^{*})\, b_1 \wedge \cdots \wedge b_n\bigr),
\end{gather}
ce qui établit la deuxième formule (11).

En particulier, supposons $E = F$, alors on voit \uncertain{par} l'identification $u \to \Lambda u$ de $L(E,F)$ dans $L(\Lambda^{(2)} E, \Lambda^{(2)} F)$ \uncertain{est} \struck{multiplicative (formule (10))} \struck{\ill{}} \struck{\uncertain{avec} l'involution $*$ (formule 11)} ; \struck{De plus} \ill{} transforme un opérateur \uncertain{positif} $h$ en un opérateur positif, et \ill{}, pour \struck{$\Lambda$} \ill{} $\alpha > 0$ \ill{}
\[
  \text{(12)}\qquad (\Lambda h)^{\alpha} = \Lambda(h^{\alpha}) \qquad (h \in L(E),\; h \geq 0)
\]
\note{le numéro (12) cerclé au crayon rouge}
En effet, si $h$ est positif, on a $h = k^{2}$, où $k$ est positif, donc $(\Lambda h) = \Lambda k^{2} = (\Lambda k)^{2}$, \struck{\ill{}} et $\Lambda k$ est \ill{} (car, \ill{} $\Lambda u$ \ill{}
\note{la phrase se poursuit page 65}

\page{65}
\uncertain{hermitien} \uncertain{positif}, \ill{} \uncertain{que} $\Lambda k$ \uncertain{est}, \struck{comme} \uncertain{aussi} \struck{l'} \ill{} hermitien, et hermitien positif. Enfin, la formule \uncertain{(12)}, \uncertain{n'est} \ill{} \uncertain{que} \uncertain{si} $\alpha$ \uncertain{est} \uncertain{rationnel} : \ill{} si $\alpha = 1/n$ \ill{} \uncertain{on} \uncertain{aura} $(\Lambda h^{1/n})^{n} = \Lambda((h^{1/n})^{n}) = \Lambda h$, et \struck{\ill{}} \ill{} $\Lambda h^{1/n}$ \ill{} hermitien positif, il en résulte \struck{que (12)} \ill{} \uncertain{vrai} \uncertain{pour} \ill{} $(\Lambda h)^{1/n} = \Lambda(h^{1/n})$. On en conclut aussitôt que (12) est vraie pour $\alpha$ rationnel \uncertain{positif} \ill{} \uncertain{puis} \ill{} $\sqrt{2}$, \ill{} \uncertain{vraie} \uncertain{pour} $\alpha$ \uncertain{positif} \uncertain{dans} \ill{}, si $H$ est un opérateur hermitien positif dans un espace de Hilbert, \ill{} les $H^{\alpha}$ \ill{} \uncertain{continûment} \uncertain{de} $\alpha > 0$, donc l'application $u \to \Lambda u$ de $L(E)$ dans $L(\Lambda^{(2)} E)$ \uncertain{étant} \ill{} continue, on voit que les fonctions $(\Lambda h)^{\alpha}$ et $\Lambda(h^{\alpha})$ sont fonctions continues de $\alpha$, qui sont égales pour $\alpha$ rationnel $> 0$, \ill{}, \uncertain{donc} égales pour tout $\alpha > 0$, ce qui achève d'établir (12).

Résumons l'essentiel de ces résultats :

\emph{Proposition 1.} \note{trait vertical en marge gauche le long de l'énoncé} Soient $E$, $F$ deux espaces de Hilbert, soient $u_i$ ($i = 1, \ldots, n$) des applications linéaires continues de $E$ dans $F$, alors $u_1 \wedge \cdots \wedge u_n$ est \struck{continue} \uncertain{prolongeable} \uncertain{en} \uncertain{une} \uncertain{application} \uncertain{linéaire} \uncertain{continue} de $\Lambda^{n\,(2)} E$ dans $\Lambda^{n\,(2)} F$ \note{« $\Lambda^{n\,(2)}$ » ajouté en interligne}, \struck{satisfaisant} de norme $\leq \|u_1\| \cdots \|u_n\|$. L'identification $u \to \Lambda u$ de $L(E)$ dans $L(\Lambda^{(2)} E)$ est multiplicative, permute à l'involution $*$, transforme opérateurs positifs en opérateurs positifs, et
\note{la phrase se poursuit page 66, sous la ligne de tête}

\page{66}
\marginal{\uncertain{à} \uncertain{mettre} \uncertain{après} \struck{\ill{}} \uncertain{Weyl} \uncertain{sur} \uncertain{les} \uncertain{propriétés} \uncertain{des} \ill{} \uncertain{hilbertiens}} \note{ligne de tête, avec un mot au crayon rouge, peut-être « À faire »}

\emph{Proposition 2.} \struck{Supposons} Soit $u$ un opérateur compact dans \ill{} l'espace de Hilbert $E$, $(\lambda_i)$ la suite de ses valeurs \struck{propres} \note{en interligne : « caractéristiques »}, alors \struck{les valeurs} $\Lambda u$ \ill{}, les valeurs \struck{propres} \note{en interligne : « caractères »} de $\Lambda^{(2)} E$ \struck{\ill{}} \ill{} sont les \ill{} $\lambda_{i_1} \cdots \lambda_{i_n}$ avec $i_1 < i_2 < \cdots < i_n$.

Comme l'identification $u \to \Lambda u$ \struck{\ill{}} de $L(E)$ \struck{\ill{}} dans $L(\Lambda^{(2)} E)$ est continue, \struck{si $u$ est \ill{}} \uncertain{et} \ill{}, \struck{fini} \struck{d'opérateurs de rang fini} \struck{\ill{}} ; \uncertain{il} \uncertain{suffit} \uncertain{de} \ill{} \uncertain{si} \ill{} \uncertain{d'opérateurs} \uncertain{de} \uncertain{rang} \uncertain{fini}, \ill{} \uncertain{limite} \uncertain{d'opérateurs} de rang fini \ill{} ; \struck{Alors} \uncertain{d'autre} \uncertain{part} si $u$ est de rang fini, alors $\Lambda u$ est \ill{}, \ill{} \uncertain{fini}, \ill{} limite d'opérateurs de rang fini, \uncertain{donc} \ill{} \struck{\ill{}} \note{bloc de quatre lignes biffé en boucles, où se lisent « base hilbertienne convenable », « sous forme d'une », « considérons », « un \uncertain{fini} seulement »}

Soit $E_1$ \uncertain{l'orthogonal} \uncertain{de} $u(E)$, \uncertain{et} $E_0$ \uncertain{l'orthogonal} \uncertain{de} $E_1$, \uncertain{et} $u(E) \subset E_0$. \note{lecture incertaine de la répartition des rôles de $E_0$ et $E_1$ ; la suite les fixe : $u$ s'annule sur $E_1$, $E_0$ est de dimension finie} \ill{} Soit $E_1$ \ill{} de $u(E)$, \ill{} ; alors $u$ s'annule sur \ill{} \note{souligné au crayon rouge} $E_0$ \ill{} \note{souligné au crayon rouge} \ill{} $\Lambda E$ est \ill{} à la \uncertain{somme} \ill{}
\[
  \sum_{0 \leq p \leq n} \Lambda^{p} E_0 \otimes \Lambda^{n-p} E_1
\]
\note{« $n-p$ » en interligne} \struck{$\Lambda E$} \ill{} $\Lambda E$ \uncertain{est} \ill{}, d'ailleurs pour $\Lambda^{p} E_0$ et $\Lambda^{q} E_1$, \uncertain{de} leur structure \ill{}, d'où \ill{} sur $\Lambda^{p} E_0 \otimes \Lambda^{q} E_1$, \uncertain{préhilbertienne} \uncertain{naturelle} \ill{}
\[
  \Lambda E \cong \sum_{p+q=n} \Lambda^{p} E_0 \otimes \Lambda^{q} E_1
\]
\marginal{\ill{} dans les \uncertain{divers} \ill{} ; $u$ compact dans un espace de Hilbert, \uncertain{ses} \uncertain{valeurs} \uncertain{propres} \uncertain{comptées} \uncertain{avec} \uncertain{leur} \uncertain{multiplicité} \ill{} on voit \ill{} ; à \ill{} des \ill{} \uncertain{induites} \ill{} \uncertain{propres} \ill{} \uncertain{modulaires} \ill{} valeurs propres de $(u^{*}u)^{1/2}$ \ill{} $(uu^{*})^{1/2}$. \ill{} \uncertain{Conséquences} \ill{} Pour les \ill{}, les \ill{} complexes, \ill{} $(\lambda_i(u))$ \ill{}, de façon \ill{} des valeurs \ill{} $\rho_i(u) = \bigl(\lambda_i(uu^{*}u)^{1/2}\bigr)$ \note{ainsi sur la page ; on attend $\rho_i(u) = \lambda_i((u^{*}u)^{1/2})$} \ill{} Cf. \ill{}} \note{longue note dans la marge gauche, en petite écriture, biffée de traits obliques ; c'est un premier état de la Proposition 2 et de ses conséquences}

\page{67}
\note{page très rapide, dont seules les formules et quelques membres de phrase se laissent lire}
\uncertain{comme} \uncertain{les} \uncertain{structures} \uncertain{préhilbertiennes} \ill{}, \uncertain{de} \ill{} \uncertain{de} $E$ \uncertain{et} \ill{} \uncertain{orthogonaux} \ill{} $E$, \struck{\ill{}} \ill{} $\Lambda^{(2)} E \cong \sum_{p+q=n} \Lambda^{p\,(2)} E_0 \otimes \Lambda^{q\,(2)} E_1$ \ill{} \note{lourdement biffé et réécrit} Si $u$ est \uncertain{un} \ill{} dans $E$, \ill{} \struck{$E_1$} \ill{} $(\Lambda u)(x_1 \wedge \cdots \wedge x_n) = u x_1 \wedge \cdots \wedge u x_n = 0$ \ill{} \uncertain{puisque} \ill{}, \ill{} $E_1$ \ill{}, donc $\Lambda u$ s'annule sur \struck{\ill{}} \ill{} \uncertain{si} $p \neq 0$ \ill{} $\sum_p \Lambda^{p} E_0 \otimes \Lambda^{n-p} E_1$ \ill{} $(p \neq 0)$ \ill{} \uncertain{alors} $\Lambda u$ \ill{} d'ailleurs \ill{}, il \ill{} $\Lambda^{(2)} E$ \ill{} et \ill{} $\Lambda^{n\,(2)} E_0 \otimes \Lambda^{0\,(2)} E_1$ \ill{}, on voit \ill{} $p \neq 0$. \ill{} \uncertain{l'ensemble} \ill{} $E_1$, \ill{} $\Lambda^{(2)} E_0$, et \ill{}, il \uncertain{est} \uncertain{évident} \ill{} $E_0$. Donc, ($u(E_1) = 0$, $u(E_0) \subset E_0$, et $E_0$ de dim finie) \ill{} \uncertain{on} \uncertain{peut} \ill{} \struck{$\Lambda u$ \ill{}} \ill{} $\Lambda^{n\,(2)} E$ \ill{} et les valeurs propres \ill{} $\Lambda^{n\,(2)} E_0 = \Lambda E_0$. \ill{} \struck{\ill{}} \ill{} les valeurs propres \uncertain{de} \ill{} $E_0$, \ill{}, les valeurs propres \ill{} \uncertain{dans} \ill{} \uncertain{les} \ill{}, \uncertain{l'algèbre} \uncertain{linéaire} \ill{}, les \ill{}, \ill{} valeurs propres \uncertain{distinctes}, (on \uncertain{le} \uncertain{voit} \uncertain{p.~ex.} \uncertain{en} \ill{} $u$ \ill{}, \ill{} \uncertain{matrice} \uncertain{triangulaire}), et \ill{}

\page{68}
\note{le corps de la page est en grande partie biffé ; un état corrigé du Corollaire 1 et de ses formules (21) et (22) est écrit de côté dans la marge gauche, et transcrit après le corps}
\struck{\ill{}} \note{trois lignes biffées} \emph{Corollaire 1.} Soit $u$ un opérateur compact \ill{} $E$ \ill{} \struck{\ill{}} \note{bloc de cinq lignes biffé en boucles}
\[
  \text{(23)}\qquad |\lambda_1(u) \cdots \lambda_n(u)| \leq \rho_1(u) \cdots \rho_n(u)
\]
\note{formule encadrée au crayon rouge ; les numéros (21) et (22) sont ceux des formules de la marge}
En effet, d'après (21) \ill{} \struck{\ill{}} $|\lambda_1(\Lambda u)|$ \ill{} \struck{$\|\Lambda u\| = \|((\Lambda u)^{*}(\Lambda u))^{1/2}\|$} \ill{} $(u^{*}u)^{1/2}$ \ill{} $(\rho_1 \cdots)$ \ill{}, (20), \ill{} \uncertain{par} \ill{} $\|\Lambda u\|$ \ill{} ; \uncertain{deuxième} \ill{} de (23), \ill{} de (22), \ill{} $\|\Lambda h\|$ \struck{\ill{}} $\Lambda h$ est hermitien \ill{}, \ill{} $\|\Lambda h\|$ \ill{} $\Lambda h$, donc \ill{}

\emph{Corollaire \uncertain{3}.} \note{aucun « Corollaire 2 » n'est visible entre les deux} Soient $E$, $F$, $G$ trois espaces de Hilbert, \ill{}
\note{la phrase se poursuit page 69}

\marginal{\emph{Corollaire 1.} Soit $u$ un opérateur compact \ill{},
\[
  \text{(21)}\qquad \lambda_1(u) \cdots \lambda_n(u) = \lambda_1(\Lambda^{n} u)
\]
\[
  \text{(22)}\qquad \rho_1(u) \cdots \rho_n(u) = \|\Lambda^{n} u\|
\]
(21) \ill{} ; (22) \ill{} $\rho_1(\Lambda u)$ \ill{} $(u^{*}u)^{1/2}$ \ill{} $= \|\Lambda u\|$ \ill{} $((u^{*}u)^{1/2})$ \ill{} ; \ill{} (23) \ill{}} \note{les numéros (21) et (22) sont cerclés, les deux formules encadrées au crayon rouge}

\page{69}
\struck{espace de Hilbert $H$,} \note{accolade en marge gauche le long des cinq premières lignes} \uncertain{Soient} \ill{} $u$ \uncertain{un} \uncertain{opérateur} \ill{} de $E$ dans $F$, $v$ \uncertain{un} \uncertain{opérateur} \ill{} de $F$ dans $G$, \ill{} $w = vu$ \note{lecture incertaine ; la formule (24) est écrite pour $uv$}. \struck{Soit} $(\rho_i(w))$ \uncertain{la} \ill{} valeurs propres de $\sqrt{w^{*}w}$. Alors, on a
\[
  \text{(24)}\qquad \rho_1(uv) \cdots \rho_n(uv) \leq \bigl(\rho_1(u)\, \rho_1(v)\bigr) \cdots \bigl(\rho_n(u)\, \rho_n(v)\bigr)
\]
\note{formule encadrée au crayon rouge ; « $\rho_1(uv) \cdots \rho_n(uv)$ » souligné}
En effet, c'est évident pour $n = 1$, car on a $\rho_1(w) = \|(w^{*}w)^{1/2}\| = \|w\|$, \struck{\ill{}} et $\|uv\| \leq \|u\|\, \|v\|$. \uncertain{Or} \ill{}, \note{bloc entre crochets} \ill{} \struck{\ill{}} $\rho_i(uv)^{2} = \rho_i((uv)^{*}(uv))$ \ill{}
\[
  \bigl(\rho_1(uv) \cdots \rho_n(uv)\bigr)^{2} = \lambda_1\bigl((uv)(uv)^{*}\bigr) \cdots \lambda_n\bigl((uv)(uv)^{*}\bigr) = \lambda_1\bigl(\Lambda((uv)(uv)^{*})\bigr) \leq \ill{}
\]
\ill{} \uncertain{de} (23), \ill{} \struck{\ill{}} $\|\Lambda(uv)\|$, i.e. $\|(\Lambda u)(\Lambda v)\|$, \struck{\ill{}} $\leq \|\Lambda u\|\, \|\Lambda v\|$, \ill{} \uncertain{en} \uncertain{vertu} \uncertain{de} \ill{} de (23), \ill{} (3) et (4), \ill{} \uncertain{que} \ill{} \uncertain{suivant}.

\uncertain{L'idée} \uncertain{de} \uncertain{la} \uncertain{démonstration} \ill{} \uncertain{est} \uncertain{due} \uncertain{à} H.~Weyl, \uncertain{ainsi} \uncertain{que} \ill{} \uncertain{inégalités} \uncertain{du} \ill{} \uncertain{suivant}.

\struck{\ill{}} \note{le bas de la page porte au crayon un bloc d'une dizaine de lignes, effacé ou très pâle, qui s'achève sur « Corollaire 3 » ; illisible}

\section{Puissances extérieures d'op. de Fredholm, inégalités diverses}

\note{Son titre, en tête de la page 70 : « 3. Puissances extérieures d'op. de Fredholm, inégalités diverses. » Les pages 70 à 73 sont une mise au net, sur des feuilles plus blanches déchirées en haut ; la numérotation des formules repart de (1). Les pages 76 à 79 en sont un état antérieur, transcrit à sa place.}

\page{70}
\emph{Prop.~1.}
\[
  \text{(1)}\qquad \|u_1 \wedge \cdots \wedge u_n\|_1 \leq \frac{1}{n!}\, \|u_1\|_1 \cdots \|u_n\|_1
\]

\emph{Corollaire 1.}
\[
  \text{(2)}\qquad \alpha_n(u_1, \ldots, u_n) = \operatorname{Tr}(u_1 \wedge \cdots \wedge u_n)
\]
\[
  \text{(3)}\qquad \alpha_n(u) = \operatorname{Tr} \Lambda^{n}(u)
\]

\emph{Corollaire 2.}
\[
  \text{(4)}\qquad |\alpha_n(u_1, \ldots, u_n)| \leq \frac{1}{n!}\, \|u_1\|_1 \cdots \|u_n\|_1 \qquad \text{---}\qquad |\alpha_n(u)| \leq \frac{1}{n!}\, \|u\|_1^{n}
\]
\note{les deux formules sont enfermées dans une boucle ; au-dessous, biffé : « \struck{(Fredholm \ill{}) $\leq \frac{1}{n!} \|u\|_1^{n}$} »}

\struck{\emph{Corollaire 3.} (6) $\|\Lambda u\|_1 = \alpha_n(|u|)$ \quad (7) $\|\Lambda u\|_2 = \sqrt{\alpha_n(u^{*}u)}$} \note{biffé d'un zigzag}

\struck{\emph{Corollaire 4.}} (8) $|\alpha_n(u)| \leq \alpha_n(|u|)$ \note{à droite, dans un cadre : « ($\struck{|\alpha_n(u)| \leq}$ $\|\Lambda u\|_1 = \alpha_n(|u|) \leq \frac{1}{n!} \|u\|_1^{n}$) »}

\emph{Corollaire 3.}
\[
  \text{(5)}\qquad |\alpha_n(u)| \leq \struck{\ldots}\; \alpha_n(|u|) = \|\Lambda u\|_1 \leq \frac{1}{n!}\, \|u\|_1^{n}
\]
\note{formule encadrée ; les numéros (6) à (8) ci-dessus sont ceux du premier état, le Corollaire 3 définitif reprenant le numéro (5)}
Car on a $|\alpha_n(u)| \leq \|\Lambda u\|_1$ en vertu de (3), et $\alpha_n(|u|) = \operatorname{Tr}(\Lambda |u|) = \operatorname{Tr} |\Lambda u| = \|\Lambda u\|_1$ (\uncertain{dans} \ill{} \uncertain{définition}), enfin $\|\Lambda u\|_1 \leq \frac{1}{n!} \|u\|_1^{n}$ (formule (1)).

\struck{? \emph{Corollaire 4.} (\ill{}) $\|\Lambda u\|_1 = \alpha_n(|u|)$, $\|\Lambda u\|_2 = \sqrt{\alpha_n(u^{*}u)}$} \note{biffé}

\emph{Corollaire 4.}
\[
  \text{(6)}\qquad |\det(1+u)| \leq \det(1+|u|)
\]

\emph{Théorème 1.} Soit $u \in L^{1}(E)$, alors \struck{$\det(1+zu)$} $(\lambda_i(u)) \in l^{1}$, et $\det(1+zu)$ \uncertain{est} de genre $0$ :
\[
  \text{(7)}\qquad \det(1+zu) = \prod_i \bigl(1 + z\,\lambda_i(u)\bigr)
\]
Équivalent à
\[
  \text{(8)}\qquad \alpha_n(u) = \sum_{i_1 < \cdots < i_n} \lambda_{i_1}(u) \cdots \lambda_{i_n}(u) \qquad (\uncertain{\text{première}} \ill{} \Longrightarrow
\]
\[
  \text{(9)}\qquad \operatorname{Tr} u = \sum_i \lambda_i(u)
\]
d'où résulte de nouveau (8) en l'appliquant à $\Lambda u$.

\uncertain{Simplifions}
\[
  \text{(10)}\qquad \alpha_n(|u|) = \|\Lambda u\|_1 = \sum_{i_1 < \cdots < i_n} \rho_{i_1}(u) \cdots \rho_{i_n}(u)
\]
Pour la \ill{}, on est ramené au cas où $u$ est hermitien positif, où c'est facile.

\emph{Prop 2.} \struck{$\|(\Lambda^{p} u) \wedge (\Lambda^{q} v)\|_1 \leq \|\Lambda^{p} u\|_1\, \|\Lambda^{q} v\|_1$}
\[
  \text{(11)}\qquad \|(\Lambda^{p} u) \wedge (\Lambda^{q} v)\|_1 \leq \frac{p!\, q!}{(p+q)!}\, \|\Lambda^{p} u\|_1\, \|\Lambda^{q} v\|_1 = \frac{p!\, q!}{(p+q)!}\, \alpha_p(|u|)\, \alpha_q(|v|)
\]
(Le second membre \ill{} à deux \uncertain{facteurs}, \uncertain{l'écriture} $(\Lambda^{p} u) \wedge (\Lambda^{q} v)$ \ill{} \uncertain{une} \uncertain{notation} \ill{}) \note{parenthèse soulignée au crayon rouge}

\page{71}
Soit $u = \sum \rho_i\, \bar{b}_i \otimes a_i$ \note{le $b$ est surchargé sur un $a$}, $\rho_i = \rho_i(u)$, etc., $v = \sum \sigma_j\, \bar{d}_j \otimes c_j$, $\sigma_j = \sigma_j(v)$, d'où
\begin{gather}
  (\Lambda^{p} u) \wedge (\Lambda^{q} v) = \frac{1}{(p+q)!} \sum_{\substack{i_1 < \cdots < i_p \\ j_1 < \cdots < j_q}} \rho_{i_1} \cdots \rho_{i_p}\, \sigma_{j_1} \cdots \sigma_{j_q}\; \times \\
  \overline{b_{i_1} \wedge \cdots \wedge b_{i_p} \wedge d_{j_1} \wedge \cdots \wedge d_{j_q}} \otimes \bigl(a_{i_1} \wedge \cdots \wedge a_{i_p} \wedge c_{j_1} \wedge \cdots \wedge c_{j_q}\bigr)
\end{gather}
Or \ill{} \uncertain{les} $(b_i)$ \ill{} \uncertain{dans} \ill{}, \ill{}, d'où
\begin{gather}
  \|(\Lambda^{p} u) \wedge (\Lambda^{q} v)\|_1 \leq \frac{1}{(p+q)!}\, p!\, q! \sum \rho_{i_1} \cdots \rho_{i_p}\, \sigma_{j_1} \cdots \sigma_{j_q} \\
  = \frac{p!\, q!}{(p+q)!} \Bigl(\sum_{i_1 < \cdots < i_p} \rho_{i_1} \cdots \rho_{i_p}\Bigr) \Bigl(\sum_{j_1 < \cdots < j_q} \sigma_{j_1} \cdots \sigma_{j_q}\Bigr) = \frac{p!\, q!}{(p+q)!}\, \|\Lambda^{p} u\|_1\, \|\Lambda^{q} v\|_1
\end{gather}

\emph{Corollaire 1.} Si $u, v \in L^{1}(E,F)$, alors \struck{\ill{}} \uncertain{pour} $r \geq 0$
\[
  \text{(12)}\qquad \struck{\det}\; \det(1 + r|u+v|) \leq \det(1 + r|u|) \bigl(\det(1 + r|v|)\bigr)
\]
i.e. \ill{}
\[
  \text{(13)}\qquad \prod_i \bigl(1 + r\rho_i(u+v)\bigr) \leq \Bigl(\prod_i \bigl(1 + r\rho_i(u)\bigr)\Bigr) \Bigl(\prod_i \bigl(1 + r\rho_i(v)\bigr)\Bigr)
\]
On va \ill{} \uncertain{montrer} que l'inégalité \ill{}, \uncertain{donc} \uncertain{que}

\emph{Corollaire 2.} Pour tout \ill{} $n$,
\[
  \text{(14)}\qquad \alpha_n(|u+v|) \leq \sum_{0 \leq p \leq n} \alpha_p(|u|)\, \alpha_{n-p}(|v|)
\]
\note{formule encadrée}
On \ill{} $|u+v| = U(u+v)$, où $U \in L(F,E)$, $\|U\| \leq 1$, \note{en interligne : « $\alpha_n(|u+v|) = \|\Lambda^{n}(u+v)\|_1$ »} \struck{donc $\alpha_n(|u+v|) = \|\Lambda U(u+v)\|_1 = \Lambda U\, \Lambda(u+v)$}
\[
  \Lambda^{n}(u+v) = \sum_{p=0}^{n} \frac{n!}{p!\,(n-p)!}\, (\Lambda^{p} u) \wedge (\Lambda^{n-p} v), \qquad \text{d'où}
\]
\[
  \alpha_n(|u+v|)\; \|\Lambda^{n}(u+v)\|_1 \leq \sum_{p=0}^{n} \frac{n!}{p!\,(n-p)!}\, \|(\Lambda^{p} u) \wedge (\Lambda^{n-p} v)\|_1 \leq \sum \struck{\frac{n!}{p!\,(n-p)!}}\; \alpha_p(|u|)\, \alpha_{n-p}(|v|).
\]
\note{ainsi sur la page, sans signe entre $\alpha_n(|u+v|)$ et la norme ; le facteur biffé de la dernière somme est remplacé en interligne par $\frac{p!\,(n-p)!}{n!}$, dont le produit avec le coefficient binomial donne 1}

\uncertain{Prouvons} \uncertain{2} \note{deux mots centrés, comme un titre}

\emph{Théorème 2.} Pour $u, v \in L_0(E,F)$, et \note{en interligne : « $r \geq 0$ »} \struck{\ill{}} $n$, on a
\[
  \text{(15)}\qquad \prod_{1}^{n} \bigl(1 + r\rho_i(u+v)\bigr) \leq \Bigl(\prod_{1}^{n} \bigl(1 + r\rho_i(u)\bigr)\Bigr) \Bigl(\prod_{1}^{n} \bigl(1 + r\rho_i(v)\bigr)\Bigr)
\]
\note{formule encadrée}
et plus \uncertain{précisément}, \ill{}, i.e. pour \struck{\ill{}} \ill{} $0 \leq m \leq n$, \ill{} \struck{$S_n^{m}(u+v) \leq \sum_{p \leq \ldots}^{m}$}
\[
  \text{(16)}\qquad \sigma_n^{m}(u+v) \leq \sum_{p=0}^{m} \sigma_n^{p}(u)\, \sigma_n^{m-p}(v)
\]
\note{formule encadrée ; à droite : « $w \in L_0(E,F)$ et \ill{} »}
\ill{} on a posé, pour \ill{} :
\[
  \text{(17)}\qquad \sigma_n^{m}(w) = \sum_{1 \leq i_1 < \cdots < i_m \leq n} \rho_{i_1}(w) \cdots \rho_{i_m}(w)
\]
\ill{} \uncertain{on} \ill{} (16), \ill{} $n = 1$, \ill{} (18) \note{ligne coupée au bas de la feuille}
\marginal{\ill{} $\leq 0$ / à réviser} \note{de côté dans la marge gauche, avec un trait le long du Théorème 2}

\page{72}
On va prouver la proposition \note{deux traits au crayon rouge} suivante :

\emph{Prop 3.} Pour $u \struck{, v} \in L_0(E,F)$, et $m \leq n$
\[
  \text{(19)}\; \struck{\text{(18)}}\qquad \sigma_n^{m}(u) = \sup_{\substack{v_n \in L(F,E) \\ \|v_n\| \leq 1,\; \operatorname{rang} v_n \leq n}} |\alpha_m(v_n u)| = \sup_{v_n} \{\alpha_m(|v_n u|)\}
\]
\note{les deux numéros sont cerclés, (18) biffé}
\struck{\ill{}} En effet, \uncertain{posons} \ill{}, pour \ill{} l'inégalité
\[
  \struck{|\alpha_m(v_n u)| = |\operatorname{Tr} \Lambda^{m} v_n u| \leq \ill{}} \qquad \struck{\leq \|\Lambda^{m} \ill{}\| \cdot \|\Lambda^{m} u\|_1 \leq \ill{}}
\]
\note{deux lignes biffées} \ill{}, \ill{} $\rho_i(v_n u) \leq \rho_i(u)$ \ill{}, \uncertain{il} \uncertain{suffit} \uncertain{d'établir} $|\alpha_m(v_n u)| \leq \sigma_n^{m}(v_n u)$, \ill{} si $u \in L(E,E)$ est de rang $\leq n$, \ill{} on a
\[
  |\alpha_m(u)| \leq \{\alpha_m(|u|)\} \leq \sigma_n^{m}(u).
\]
\ill{} $u$ \struck{\ill{}} $\in \ill{}(E)$, et
\[
  |\alpha_m(u)| = |\alpha_m(w)| \leq \struck{\ill{}} \sum_{1 \leq i_1 < \cdots < i_m \leq n} \rho_{i_1}(w) \cdots \rho_{i_m}(w)
\]
\uncertain{et} \ill{} \struck{\ill{}} $w = \alpha^{*} u \alpha$ ($\alpha$ appl. isométrique \uncertain{$u(E)$} $\to E$) \note{lecture incertaine de la source de $\alpha$}, d'où $\rho_i(w) \leq \|\alpha^{*}\|\, \|\alpha\|\, \rho_i(u) \leq \rho_i(u)$, \uncertain{la} \uncertain{formule} \ill{}.

\struck{\ill{}} \uncertain{Montrons} \ill{} \uncertain{2}. \ill{} (19) \ill{} \note{souligné au crayon rouge} \ill{} $E$ de dimension $\geq n$, \ill{} \struck{\ill{}}, $u = \sum \rho_i\, \bar{e}_i \otimes f_i$, \ill{} \uncertain{posons} $v_n = \sum_{1}^{n} \bar{f}_i \otimes e_i$, d'où $v_n u = \sum_{1}^{n} \rho_i\, \bar{e}_i \otimes e_i$, et $\alpha_m(v_n u) = \sigma_n^{m}(u)$, (\uncertain{remarquer} \uncertain{que} si $F = E$, et $u$ hermitien, alors on peut prendre $v_n$ \uncertain{projecteur} \uncertain{hermitien} \ill{} $E$ de dimension $\leq n$). \uncertain{Prouvons} \uncertain{maintenant} (16) : \ill{} (18) \ill{} \struck{(17), \ill{} à prouver} \struck{que} pour $|\alpha_m(w(u+v))|$ \struck{$\leq \sum_{p=0} \ill{} (|wu|)$} \ill{} pour la seconde \ill{}, \ill{} $w \in L(F,E)$ telle que $\|w\| \leq 1$, $\operatorname{rang} w \leq n$, \struck{\ill{}} on \ill{} ((14))
\[
  |\alpha_m(w(u+v))| \leq \alpha_m(|wu + wv|) \leq \sum_{p=0}^{m} \alpha_p(|wu|)\, \alpha_{m-p}(|wv|)
\]
\struck{d'autre part \ill{} $\alpha_p(|wu|) = \|\Lambda^{p} wu\|_1 = \|(\Lambda^{p} w)(\Lambda^{p} u)\|_1 \leq \|\Lambda^{p} w\|\, \|\Lambda^{p} u\|_1 \leq \|\Lambda^{p} u\|_1 = \alpha_p$} \note{deux lignes biffées} \struck{et \ill{}} \ill{} (19), \ill{} $\alpha_p(|wu|) \leq \sigma_n^{p}(u)$ et $\alpha_{m-p}(|wv|) \leq \sigma_n^{m-p}(v)$, d'où la formule (18). \note{ainsi sur la page ; c'est la formule (16) qui est établie}

\page{73}
\note{seul le tiers supérieur de la page est écrit}
Le système des inégalités (17) \note{cerclé} (pour $m, n$ variables) \uncertain{et} \uncertain{la} \uncertain{constante} \ill{} \struck{\ill{}} \uncertain{pour} (15) \note{souligné au crayon rouge} \uncertain{semblent} \uncertain{les} \uncertain{inégalités} \uncertain{les} \uncertain{plus} \uncertain{fortes} \uncertain{possibles} \uncertain{majorant} \uncertain{les} \note{en interligne : « fonction des »} $\rho_i(u+v)$ en fonction des $\rho_i(u)$, $\rho_i(v)$. \uncertain{Notons} \struck{\ill{} \uncertain{l'inégalité}} (15) \ill{} \uncertain{qui} \uncertain{permettent} \uncertain{de} \ill{}, \ill{}, \uncertain{qu'elle} \ill{} \uncertain{l'inégalité} (17) \uncertain{qui} \uncertain{elle} \ill{} \uncertain{particulière}, \uncertain{pour} \struck{\ill{} le cas \ill{}} \note{plusieurs mots soulignés au crayon rouge}

\emph{Prop.~4.} Si $u \in L_0(E,F)$, \ill{} \uncertain{entier} $n \geq 0$, \ill{}
\[
  \text{(18)}\qquad \struck{\det}\; \prod_{1}^{n} \bigl(1 + \rho_i(u)\bigr) = \sup_{\substack{v_n \in L(F,E) \\ \|v_n\| \leq 1,\; \operatorname{rang} v_n \leq n}} \det(1 + v_n u)
\]
Cette \struck{démonstration} \note{en interligne : « formule »} \uncertain{est} \ill{} \uncertain{immédiate} \ill{} (17), \struck{\ill{}} \note{une ligne biffée} \struck{\ill{} \uncertain{plus} \uncertain{simple}}

\section{Sous-additivité des fonctions des valeurs singulières}

\note{Titre de l'éditeur. Deux feuilles brunes, d'une autre plume, pages 74 et 75, qui tirent de (15) des inégalités du type $\sum f(\rho_i(u+v)) \leq \sum f(\rho_i(u)) + \sum f(\rho_i(v))$.}

\page{74}
\[
  \sum_{1}^{n} \log\bigl(1 + r\rho_i(u+v)\bigr) \leq \sum_{1}^{n} \log\bigl(1 + r\rho_i(u)\bigr) + \sum_{1}^{n} \log\bigl(1 + r\rho_i(v)\bigr)
\]
\[
  S_\varphi^{n}(u+v) \leq S_\varphi^{n}(u) + S_\varphi^{n}(v), \qquad S_\varphi^{n}(w) = \sum_{1}^{n} f_\varphi\bigl(\rho_i(w)\bigr)
\]
\[
  f_\varphi(\rho) = \int_{0}^{+\infty} \frac{\varphi(r)}{\frac{1}{\rho} + r}\, dr = \int_{0}^{+\infty} \varphi(s/\rho)\, \frac{ds}{1+s}
\]
$\varphi$ fonction \uncertain{décroissante} \note{en interligne : « si $A(0) = B(0) = 0$ (\ill{}), alors »} (\uncertain{motivation} : \note{dans un cadre : « $A(r) \leq B(r)$ \ill{} $0 \leq r < +\infty$ \ill{} »} $\Longrightarrow \int \varphi(r)\, A(r)\, dr \leq \int \varphi(r)\, B(r)\, dr$ \ill{} \note{en interligne : « dans $[a,b]$ »}. \uncertain{Or}, \ill{} \uncertain{pour} \struck{\ill{}} \ill{} \note{deux mots soulignés} \ill{}, \uncertain{dans} \ill{} $\psi(r) = -\varphi'(r)$, $\varphi$ \ill{} (\struck{\ill{}} \ill{} si $A$ est dérivable
\[
  -\int \varphi'(r)\, A(r)\, dr = \bigl[-\varphi(r)\, A(r)\bigr]_{0}^{+\infty} + \int \varphi(r)\, A'(r)\, dr
\]
et la première \ill{}, d'où \ill{} $\int \varphi(r)\, A'(r)\, dr$. Si \ill{} $A$ et $B$ \ill{} \uncertain{dérivables}, alors \ill{} $\int \varphi(r)\, A'(r)\, dr \leq \int \varphi(r)\, B'(r)\, dr$ \ill{} \uncertain{pour} \ill{} $[0, +\infty[$ \ill{} fonction $\varphi$ \ill{}, \ill{} \uncertain{d'ailleurs} \ill{} et \ill{} $A'(r), B'(r) \geq 0$, \ill{}, alors \ill{} \uncertain{vrai} \uncertain{pour} \ill{} \uncertain{discontinuité} \ill{}, \uncertain{pourvu} \uncertain{que} \ill{} $\varphi = $ \note{un cadre vide} \ill{} $A(r) \leq B(r)$). \ill{} \struck{\ill{}} \ill{} \uncertain{retrouve} \ill{} Dans le cas \uncertain{actuel}, $\tilde{A}(r) = \sum_{1}^{n} \log(1 + r\rho_i(u))$, d'où \ill{} $A'(r) = \sum \frac{r}{1 + r\rho_i(u)}$ \note{ainsi sur la page ; on attend $\sum \frac{\rho_i(u)}{1 + r\rho_i(u)}$} \ill{}. Le cas le plus \uncertain{important} \uncertain{est} \ill{} pour $\varphi(r) = r^{-p}$, $0 < p < 1$. On \uncertain{trouve} alors
\note{la phrase se poursuit page 75}

\page{75}
\[
  f_p(\rho) = \rho^{p} \int \frac{s^{-p}}{1+s}\, ds = c^{te} \cdot \rho^{p}, \qquad \text{d'où, en posant}
\]
\struck{Système} \uncertain{a.}
\[
  S_p^{(n)}(u) = \sum_{1}^{n} \rho_i(u)^{p}
\]
\[
  S_p^{(n)}(u+v) \leq S_p^{(n)}(u) + S_p^{(n)}(v) \qquad \text{en particulier}
\]
\[
  S_p(u+v) \leq S_p(u) + S_p(v)
\]
\note{les deux inégalités sont encadrées}
i.e. $\sum_{1}^{n} \rho_i(u+v)^{p} \leq \sum_{1}^{n} \bigl(\rho_i(u)^{p} + \rho_i(v)^{p}\bigr)$.

\emph{Remarque.} Les \ill{}, \ill{} $p \geq 1$ \ill{} \uncertain{pas} \uncertain{dans} \ill{} ; \uncertain{le} \uncertain{cas} $p$ \ill{} Weyl ; on trouve \ill{} de \ill{}.

\emph{Question.} \uncertain{Trouver} \ill{} $f$ \ill{} fonction $f$ d'une variable positive telle que $f^{(n)}(u) = \sum_{1}^{n} f(\rho_i(u))$, \ill{}
\[
  f^{(n)}(u+v) \leq f^{(n)}(u) + f^{(n)}(v) \qquad \text{quel que soit } n.
\]
\uncertain{Il} \uncertain{suffit} \uncertain{pour} \ill{} \uncertain{que} $f^{\infty}(u+v) \leq f^{\infty}(u) + f^{\infty}(v)$ \ill{} $u \in L(E)$, \ill{}

\uncertain{Réponse} \ill{} \uncertain{arbitraire} \ill{}

\uncertain{Réf.} \ill{}

\uncertain{Généralisation} \uncertain{à} \ill{}

\section{Puissances extérieures d'op. de Fredholm : premier état}

\note{Titre de l'éditeur. Les pages 76 à 79, feuilles brunes, sont l'état antérieur de la mise au net des pages 70-73 ; il en porte la numérotation (13) à (23), et la consigne au crayon rouge qui ouvre la page 76 le renvoie à la rédaction définitive.}

\page{76}
\ill{} \note{ligne de tête coupée par la déchirure, où se lit « \uncertain{positifs} »}

\marginal{Mettre 1 : N° 2 et cor. N° ---} \note{au crayon rouge, souligné d'un trait rouge}

\emph{Proposition 3.} Soient $E$ et $F$ deux espaces de Hilbert, soient $u_i$ ($i = 1, \ldots, n$) des \uncertain{applications} \struck{\ill{}} de Fredholm de $E$ dans $F$, alors $u_1 \wedge \cdots \wedge u_n$ est une application de Fredholm de $\Lambda^{n} E$ dans $\Lambda^{n} F$, et
\[
  \text{(13)}\qquad \|u_1 \wedge \cdots \wedge u_n\|_1 \leq \frac{1}{n!}\, \|u_1\|_1 \cdots \|u_n\|_1
\]
\note{formule encadrée au crayon rouge}
En effet, il suffit de \ill{} \uncertain{vérifier} \uncertain{quand} \ill{} \struck{$u_i = a_i' \otimes$} $u_i = \bar{b}_i \otimes a_i$, \struck{\ill{}} \ill{} avec $\|b_i\| \leq 1$, $\|a_i\| \leq 1$, alors $\|u_1 \wedge \cdots \wedge u_n\|_1 \leq \frac{1}{n!}$ \ill{} \struck{\ill{}} \ill{}, \ill{} \uncertain{il} existe une \ill{} de $(\bar{E} \mathbin{\hat{\otimes}} E)$ dans $\bar{F} \mathbin{\hat{\otimes}} F$ (\uncertain{où} $F = \Lambda^{n} E$) \ill{} $(u_i)$ \ill{} qui coïncide avec $u_1 \wedge \cdots \wedge u_n$ \uncertain{pour} \uncertain{les} \ill{} \uncertain{décomposables}, \ill{} $L(F)$. Or, pour $u_i = \bar{b}_i \otimes a_i$, \ill{} (Chap.~\ill{}, N° 1, formule (8))
\[
  \struck{\ill{}}\quad u_1 \wedge \cdots \wedge u_n = \frac{1}{n!}\, (\bar{b}_1 \wedge \cdots \wedge \bar{b}_n) \otimes (a_1 \wedge \cdots \wedge a_n), \qquad \text{d'où}
\]
\[
  \|u_1 \wedge \cdots \wedge u_n\|_1 \leq \frac{1}{n!}\, \|b_1 \wedge \cdots \wedge b_n\|\, \|a_1 \wedge \cdots \wedge a_n\| \leq \frac{1}{n!} \prod \|a_i\|\, \|b_i\|
\]
(formule (5)), d'où \ill{} (13).

\emph{Corollaire 1.} Soit $E$ un espace de Hilbert, et $u_i$ ($i = 1, \ldots, n$) \struck{des opérateurs} \uncertain{des} \uncertain{opérateurs} de Fredholm dans $E$. Alors $\Lambda^{n} u$ est \uncertain{un} \uncertain{opérateur} de Fredholm dans $\Lambda^{n} E$, et on a
\[
  \text{(14)}\qquad \alpha_n(u_1, \ldots, u_n) = \operatorname{Tr}(u_1 \wedge \cdots \wedge u_n), \qquad \alpha_n(u) = \operatorname{Tr}(\Lambda^{n} u) \quad (\uncertain{\text{en particulier}})
\]
\note{les deux formules encadrées ensemble au crayon rouge}
(où $\alpha_n(u)$ est défini dans Chap.~1, \S 2, N° 1) \note{souligné au crayon rouge} En effet, (14) \uncertain{est} \uncertain{vraie} \uncertain{pour} \ill{} \uncertain{quand} \ill{}
\note{la phrase se poursuit page 77}

\page{77}
\ill{} (voir Chap.~1, \S 1, N° 2, \ill{} $(i \ldots)$), et \ill{} des \ill{} \struck{\ill{}} \ill{} \note{un membre de phrase souligné} \ill{} (13), \ill{}.

\emph{Corollaire 2.} \uncertain{Soient} $u_1, \ldots, u_n$ \ill{}
\[
  \text{(15)}\qquad |\alpha_n(u_1, \ldots, u_n)| \leq \frac{1}{n!}\, \|u_1\|_1 \cdots \|u_n\|_1
\]
\note{formule encadrée au crayon rouge}
(\uncertain{Cela} \uncertain{résulte} \uncertain{aussi} \ill{} \uncertain{inégalités} \uncertain{que} \uncertain{le} \ill{} du Chap.~1, \S 2, N° 2, \uncertain{prop}~1, \ill{} \uncertain{$\leq 1$}).

\emph{Corollaire 3.} Soit \struck{$u \in L^{1}(E)$} \ill{} \struck{\uncertain{op.} de Fredholm dans un espace de Hilbert} $E$. Alors \ill{} $u$ \ill{} \uncertain{op.} de Fredholm \uncertain{de} \ill{} $E$, \ill{}, et
\[
  \text{(16)}\qquad \|\Lambda^{n} u\|_1 = \alpha_n(|u|) = \alpha_n\bigl(\sqrt{u^{*}u}\bigr)
\]
\note{le dernier membre est écrit en interligne, souligné au crayon rouge}
Si $u$ est \ill{} \uncertain{op.} \uncertain{de} \uncertain{Hilbert-Schmidt} \ill{}, alors \ill{} \struck{\ill{}} \ill{} $\Lambda^{n} u$ \note{deux lignes biffées}, et
\[
  \text{(17)}\qquad \|\Lambda^{n} u\|_2 = \sqrt{\alpha_n(u^{*}u)}
\]
\struck{On a, si $u \in L^{1}(E)$, \ill{} formule (22)) \ill{}} \note{biffé}
$\|\Lambda^{n} u\|_1 = \operatorname{Tr}\bigl((\Lambda^{n} u)^{*}(\Lambda^{n} u)\bigr)^{1/2}$ (N° 1, \struck{formule} \ill{}) \struck{\ill{}} $\bigl((\Lambda^{n} u)^{*}(\Lambda^{n} u)\bigr)^{1/2} = \Lambda\bigl((u^{*}u)^{1/2}\bigr)$, d'où $\|\Lambda^{n} u\|_1 = \operatorname{Tr} \Lambda\bigl((u^{*}u)^{1/2}\bigr) = \alpha_n\bigl(\sqrt{u^{*}u}\bigr)$ \struck{(formule \ill{} (17))} (formule (14)). \struck{Si $u \in L^{2}(E)$, \ill{} $(\Lambda u)^{*}(\Lambda u) = \Lambda(u^{*}u)$ \ill{} $u^{*}u \in L^{1}(E)$ \ill{}} \note{bloc de trois lignes biffé en boucles} \ill{} Fredholm \ill{} (N° 1, th.~4), donc $(\Lambda^{n} u)^{*}(\Lambda^{n} u) = \Lambda^{n}(u^{*}u)$ \uncertain{est} \ill{} Fredholm, et donc $\Lambda^{n} u$ est \ill{} (N° 1, th.~3, corollaire 3) \ill{} $(\|\Lambda^{n} u\|_2)^{2} = \operatorname{Tr}(\Lambda^{n} u)^{*}(\Lambda^{n} u) = \alpha_n(u^{*}u)$ (formule (14), d'où (17))

\page{78}
\note{la moitié inférieure de la page, à partir de la démonstration du Théorème 2, est barrée de deux grands traits obliques, ainsi que toute la page 79}

\emph{Corollaire 4.} Soit $u$ un \uncertain{op.} de Fredholm dans l'espace de Hilbert $E$. Alors
\[
  \text{(18)}\qquad |\alpha_n(u)| \leq \alpha_n(|u|)
\]
\note{formule encadrée au crayon rouge}
(\uncertain{En} \uncertain{d'autres} \uncertain{termes}) \uncertain{les} \ill{} de la \uncertain{fonction} \ill{} $\det(1+zu)$ \ill{}, \ill{} Fredholm \ill{} $\det(1 + z|u|)$ \ill{} \uncertain{On} \ill{}.

\emph{Corollaire 5.} \uncertain{Pour} \ill{}
\[
  \text{(19)}\qquad |\det(1+zu)| \leq \det(1 + |z|\,|u|) \qquad \struck{(\ill{})}
\]
\uncertain{En} \uncertain{effet} : $|\alpha_n(u)| = |\operatorname{Tr} \Lambda^{n} u| \leq \|\Lambda^{n} u\|_1 = \alpha_n(|u|)$ \note{un mot souligné au crayon rouge devant la ligne} \ill{} (18), \uncertain{On} \uncertain{a} \ill{} (formule (16)).

\emph{Théorème 2.} Soit $u$ une application de Fredholm de l'espace de Hilbert $H$. Alors \uncertain{les} \ill{}, $(\lambda_i)$ \ill{}, \struck{\ill{}} \note{deux lignes biffées} on a
\[
  \text{(20)}\qquad \det(1+u) = \prod_i (1 + \lambda_i)
\]
\note{formule encadrée au crayon rouge}
\ill{} (21) ; \ill{}, et \ill{} $zu$ au lieu de $u$, \ill{} la fonction \uncertain{entière} de Fredholm $\det(1+zu)$ \ill{} de genre $0$ ;
\[
  \text{(21)}\qquad \det(1+zu) = \prod_i (1 + z\lambda_i)
\]
\emph{Démonstration.} --- \ill{} des N° \ill{}, \ill{} ; \ill{} (20) \ill{} \struck{\ill{}} \note{le reste de la page, six lignes, est biffé en boucles par-dessus les traits obliques ; on y lit « (20) », « (22) », « des N° » et « \uncertain{hypothèse} »}

\marginal{
\[
  \text{(22)}\qquad \alpha_n(u) = \sum_{i_1 < \cdots < i_n} \lambda_{i_1} \cdots \lambda_{i_n}
\]
\[
  \text{(23)}\qquad \operatorname{Tr} u = \sum \lambda_i(u)
\]
\ill{} (22') \ill{}} \note{de côté dans la marge gauche ; la formule (23) est encadrée au crayon rouge}

\page{79}
\note{toute la page est barrée de deux traits obliques ; c'est la suite de la démonstration du Théorème 2}
\uncertain{Réunissons} \ill{}, \uncertain{Il} \ill{} \uncertain{résulte} \uncertain{de} \ill{} \struck{\ill{}} \struck{$\sum \ill{} \leq \ill{}$} \ill{} \uncertain{Soient} \struck{\ill{}} $x_1, \ldots, x_n$ des vecteurs \ill{} \struck{\ill{} $x_i$, \ill{}} \ill{} i.e. \ill{}, \ill{} $\sum |z_i| \leq \|u\|_1$ \note{$z_i$ désigne ici les valeurs propres, $\lambda_i$ plus bas}, \ill{} soit en effet \ill{} $p$ \struck{\ill{}} \uncertain{le} \uncertain{projecteur} \uncertain{orthogonal} \ill{} \uncertain{engendré} \uncertain{par} \uncertain{les} $x_i$, \ill{}
\[
  |z_1 \cdots z_n| = |\operatorname{Tr} pu| \leq \|pu\|_1 \leq \|p\|\, \|u\|_1 = \|u\|_1 .
\]
\note{ainsi sur la page ; la borne $\|u\|_1$ vaut pour une somme $\sum z_i$ plutôt que pour un produit}
\ill{} \uncertain{comme} \uncertain{les} \uncertain{zéros} \ill{} de $\det(1 - zu)$ \ill{}, \ill{} \uncertain{les} \uncertain{coefficients} \uncertain{de} $\det(1 - zu)$ \uncertain{sont} \ill{} (prop.~2, corollaire 4). \uncertain{Nous} \ill{} \uncertain{fonction} \uncertain{entière} \ill{} \uncertain{les} \ill{}, et \ill{} \uncertain{majorés} \uncertain{en} \uncertain{modules} \uncertain{par} \uncertain{ceux} \ill{} \uncertain{genre} $0$ ; \ill{} \struck{\ill{}} $0$. \struck{Pour \ill{} une \ill{} fonction} \ill{} \uncertain{que} $\det(1 + zh)$ \ill{} \uncertain{de} \uncertain{genre} \ill{} \struck{\ill{} (22) \ill{}} ; \struck{Mais \ill{} l'identité (22) $\alpha_n(u) = \sum_{i_1 < \cdots < i_n} z_{i_1} \cdots z_{i_n}$ \ill{} résulte immédiatement de l'expression} \note{deux lignes biffées, avec la formule (22) cerclée} \ill{} $(e_i)$ \ill{} \uncertain{un} \uncertain{système} \uncertain{orthonormal} \ill{}, $u = \sum \lambda_i\, \bar{e}_i \otimes e_i$ (\ill{}) \struck{En effet, \ill{}} \uncertain{d'où} \uncertain{donc}
\[
  \alpha_n(u) = \sum_{i_1 \ldots i_n} \struck{\alpha_n(z_i e \ldots)}\; \lambda_{i_1} \cdots \lambda_{i_n}\; \alpha_n\bigl(\bar{e}_{i_1} \otimes e_{i_1}, \ldots, \bar{e}_{i_n} \otimes e_{i_n}\bigr)
\]
\ill{} \struck{\ill{}} \ill{}, \ill{} ; \ill{} $\frac{1}{n!}$ ; \ill{} \uncertain{les} $i_1, \ldots, i_n$ \ill{} \uncertain{distincts}, \ill{} (\ill{}). \ill{} \uncertain{différent} \ill{} (23).
\marginal{\ill{} $\sum |\lambda_i| \leq \sum$ \ill{} $u$ \ill{}} \note{de côté dans la marge gauche, très rapide}

\section{Fonctions convexes : hyperplans d'appui}

\note{Titre de l'éditeur. Une feuille isolée, page 80, sans rapport visible avec les feuillets précédents sinon la « Question » de la page 75.}

\page{80}
$f$ fonction convexe sur \uncertain{l'ouvert} \uncertain{convexe} $U$,
\[
  (x,y) \in G_f \iff x \in U \text{ et } y \geq f(x)
\]
Soit \struck{$x_i, y_i$} $(a,b)$ avec $b = f(a)$ ($a \in U$), \struck{donc} \uncertain{soit} $H$ \uncertain{un} \uncertain{hyperplan} \ill{} \uncertain{passant} \uncertain{par} $(a,b)$ \uncertain{ne} \uncertain{rencontrant} \uncertain{pas} \ill{} $G_f$, alors : 1) \uncertain{il} \ill{} \uncertain{a} \ill{} \struck{$y - b = \ill{}$} \note{lourdement biffé} \ill{} $\sum c_i (x_i - a_i)$ (\ill{} \uncertain{droite} $x = a$), soit $y = L_a(x)$ \ill{} ; 2) \ill{} $f(x) \geq L(x)$ pour \struck{\ill{}} $x \in U$ (\ill{} $H$, \ill{} \struck{\ill{}} ; \ill{} $H$ \uncertain{contient} $G_f$, et \ill{} \uncertain{que} \uncertain{alors} \ill{} \uncertain{pour} $y \geq L(x)$. \note{trait ondulé sous « \uncertain{détermination} »} \uncertain{détermination} \uncertain{pour} $H$ \ill{} \uncertain{pour} \ill{} \uncertain{que} \ill{} \uncertain{alors} \ill{} \uncertain{2} \ill{}), 3) \ill{} $\in E$ \ill{} ($L$ \ill{}), et \ill{} \struck{\ill{}} $a$, \ill{}
\[
  \text{4)}\qquad f(x) = \sup_a L_a(x) \qquad (\ill{})
\]
\ill{} \uncertain{si} \ill{}, \uncertain{les} $c_i$ \ill{} $\geq 0$. \uncertain{En} \uncertain{effet}, \ill{} \uncertain{croissante} \ill{} $f(x_1, a_2, \ldots, a_n)$ \ill{} et \uncertain{de} \ill{} \uncertain{première} \ill{} \uncertain{linéaire} \ill{} \uncertain{Lorsque} \ill{} \ill{} \uncertain{supposons} \ill{} $c_1 < 0$, alors \ill{} \uncertain{voisinage} \ill{} \uncertain{pour} $x_1 = a_1$, \ill{} $L(x_1, a_2, \ldots, a_n)$, \ill{} \uncertain{et} \uncertain{on} \uncertain{aurait} \ill{} $L = $ \ill{} $f$. \ill{} \uncertain{pour} \ill{} \uncertain{diminuer} \ill{}
\[
  f(x_1) \leq f(a) \leq f(a_1)
\]
\note{ainsi sur la page : les arguments sont abrégés à leur première coordonnée}
\uncertain{voisin} \uncertain{de} $a_1$.

\end{document}
